%%%%%%%%%%%%%%%%%%%%%%% CHAPTER - 2 %%%%%%%%%%%%%%%%%%%%\\
\chapter{Literature Review}
\label{C2} %%%%%%%%%%%%%%%%%%%%%%%%%%%%
\graphicspath{{Figures/Chapter-2figs/PDF/}{Figures/Chapter-2figs/}}
% \noindent\rule{\linewidth}{2pt}
%%%%%%%%%%%%%%%%%%%%%%%%%%%%%%%%%%%%%%%%%%%%%%%%%%%%%%%%%%%%%%%%%%%%%%%%%%%%%%%%%%%%%%%%%%%%%%%%%%%%%%%%%%%%%%%%%%%%%%%%%%%%%%%%%%%%

The complexity of human movement in sports has driven growing research in posture detection 
	and correction. Earlier methods based on manual observation and sensors have evolved into AI 
	and deep learning approaches, enabling more accurate and real-time exercise monitoring.

Musa Çankaya and Fatma Nur Takı~\cite{cankaya2024posture} conducted a comparative study on postural assessment and awareness in individuals receiving posture training through a digital AI-based system. 
Their research highlighted that traditional posture correction methods often fail to provide 
consistent results, while digital AI systems offer personalized and adaptive training feedback. 
The study demonstrated that AI-based posture detection helps participants maintain correct 
body alignment during daily activities. They found that AI-driven feedback mechanisms 
significantly improved posture awareness among individuals who underwent training. 
Furthermore, the system reduced the risk of musculoskeletal disorders by correcting harmful 
sitting and standing postures early. The findings strongly supported the idea that integrating AI 
into posture correction has the potential to replace conventional methods. Their conclusion 
emphasized that AI- based posture monitoring tools are not only more efficient but also 
scalable for use in occupational health and rehabilitation. 

Engdahl, C et al. ~\cite{engdahl2023dance} proposed a system to enhance the teaching of creative aspects of dance in Teacher Education (PETE). Their work emphasized that fostering creativity is a globally 
recognized educational objective and identified expressive dance as an effective medium to 
develop students’ creativity. However, they observed that traditional dance education in 
schools and PETE often focuses on reproducing pre- arranged movement sequences and 
achieving predetermined outcomes, which limits opportunities for creative exploration and 
problem-solving. Using the theoretical framework of and, specifically the concepts of smooth 
and striated spaces and experimentation, Engdahl et al. analysed interviews with eight PETE 
dance educators from different institutions in. Their findings showed that while creativity is 
mentioned in PETE curricula, it remains marginal in practice due to the dominant emphasis on 
technical correctness and physical performance in. They concluded that teacher educators play 
a crucial role in shaping how future PE teachers incorporate creativity into their teaching, and 
suggested that encouraging exploratory tasks, improvisation, problem-solving, and the use of 
music are effective strategies to promote creative movement learning in expressive dance 
education.

Fan et al. ~\cite{fan2015pose}  proposed a dual-source deep neural network that combines local appearance and holistic view for human pose estimation. Their framework was designed to handle the 
limitations of single-source models, which often struggled with occlusion and complex 
backgrounds. The integration of local body part features with global structural context 
significantly improved recognition accuracy. The model was tested in challenging 
environments, where it consistently outperformed existing benchmarks. By fusing local and 
holistic information, the network could better distinguish between similar poses with minimal 
errors. This dual-source approach proved particularly effective in sports and fitness 
applications, where body movement analysis requires high precision. The results indicated that 
deep learning frameworks could revolutionize the field of human pose estimation by ensuring 
robustness across varying conditions. The authors concluded that their framework laid the 
foundation for future improvements in pose estimation technology.
From our analysis, we found that there is no integrated automated system that effectively 
combines real-time cricket performance analysis, training feedback, and posture correction 
using Artificial Intelligence in an efficient and accessible manner.

Felzenszwalb and Huttenlocher ~\cite{felzenszwalb2005pictorial} introduced pictorial structures for object recognition, which laid a foundational framework for modern posture estimation methods. Their approach  
represented objects as collections of parts connected by spatial relationships, making it one of 
the earliest structured models in computer vision. This methodology proved highly effective 
for representing human body structures in posture recognition. Although their work was not 
originally intended for sports or exercise applications, it provided the theoretical groundwork 
for subsequent AI-based pose estimation research. Their model also influenced the 
development of part-based human detection systems that became critical in posture monitoring. 
The study showed that structured approaches are highly versatile, as they balance local 
accuracy with global structural consistency. The impact of their research continues to be visible 
in recent AI-driven posture detection frameworks.

Qin Z et al. ~\cite{qin2022gnn} proposed a novel method for Skeleton-based action recognition by 
introducing higher-order angular features into Graph neural networks (GNNs) to improve 
recognition accuracy. They noted that skeleton-based action recognition is more robust to 
background interference and easier to process, and that recent advances in GNNs have 
significantly boosted its performance by modeling skeletons as graphs— where joints are nodes 
and bones are edges. Earlier methods used only first-order joint coordinate features, and later 
second-order bone vector features, which improved performance but still struggled to 
distinguish actions with similar motion trajectories. To overcome this limitation, Qin et al. 
proposed an angular encoding (AGE) approach that captures relative movements between body 
parts through angles formed by three related joints, providing invariance to different body sizes 
and motion speeds. They demonstrated that integrating angular features into existing models 
such as Spatio- temporal graph convolutional network (ST-GCN) and decoupling GCN 
improves the classification of confusingly similar actions. Experimental results showed that 
their method achieved state-of-the-art accuracy on benchmark datasets like NTU60 and 
NTU120.

Supriya Shrivastav et al.~\cite{shrivastav2025exercise} developed an AI-based exercise posture recognition system using an estimation framework. Their work focused on real-time analysis, enabling immediate 
detection and correction of incorrect exercise postures. The system was capable of identifying 
a wide range of common exercises, making it suitable for fitness coaching and rehabilitation. 
The researchers highlighted the importance of accurate posture detection in preventing injuries 
during workouts. Their model demonstrated a high recognition rate and adaptability across 
diverse body types and exercise forms. By integrating AI-based feedback, users were able to 
refine their techniques without the need for constant trainer supervision. The study showed that 
the framework not only enhanced training efficiency but also promoted safer workout practices. 
Their results emphasized that such systems have strong potential in the fitness industry, 
especially for at-home and remote exercise programs.

Tang T and Hyun-Joo M ~\cite{tang2022sportsdance} proposed a sports dance movement detection system based on 
Pose recognition to enhance the learning and analysis of human actions and behaviours. They 
emphasized that pose recognition can standardize movements during dance teaching and 
preserve vital motion data, unlike traditional training methods that rely on repeated video 
viewing or coach feedback. Their system follows three main stages data acquisition and 
preprocessing, feature extraction and construction, and action recognition and addresses the 
low accuracy of existing methods by selecting key skeletal points and using a residual block
based detection method. It captures complete dance movement data through posture 
recognition and bone tracking, compares it with standard actions to detect deviations, and helps 
learners correct nonstandard postures. The system also supports flexible, location-independent 
training, promoting the digitalization of Sports Dance education. 

Tian et al. \cite{tian2024bigdata} explored the integration of Big Data and Artificial Intelligence in football to enhance training, competition, and refereeing. The study highlighted how data analytics and 
AI-driven tools such as VAR, GoalControl, and Scout Advisor improve tactical decision
making, player evaluation, and match fairness. The authors also emphasized the role of VR, 
AR, and 3D reconstruction in creating immersive training environments. However, challenges 
such as high implementation costs, privacy concerns, and lack of standardization were noted. 
The paper concluded that combining Big Data and AI provides a scientific and intelligent 
framework for modern football development, enabling data-driven coaching, injury prevention, 
and real-time performance optimization.

Wang and Zhang~\cite{wang2024posture} explored AI-powered athletic posture detection in sports training through empirical research. They argued that traditional coaching methods lack precision in monitoring 
posture and movements, especially in high-speed athletic activities. By incorporating AI 
algorithms, their system was able to deliver real-time feedback to athletes, allowing for  
immediate corrections. Their study revealed that the AI framework significantly enhanced the 
efficiency of training by reducing repetitive mistakes in postures. Moreover, the system could 
detect subtle deviations that were often overlooked by human trainers. This led to better 
performance outcomes, as athletes could fine-tune their techniques with data-driven insights. 
The results highlighted that AI-based monitoring is crucial for modern sports training, where 
even small posture errors can impact performance. Ultimately, their findings stressed the 
importance of merging technology with sports science for achieving peak athletic performance. 

Yang, X~\cite{yang2022dance} proposed a dance teaching system based on to enhance the evaluation and 
prediction of students’ dance performance. In this study, a special performance prediction 
model was developed where the physical form quality index was considered as the independent 
variable and the special achievement was taken as the dependent variable. Through association 
analysis and cluster analysis, six specific items with higher sensitivity related to body shape, 
physical quality, and special quality were identified. A model was then adopted to build the 
prediction framework, and its reliability and error were analysed and compared with a 
traditional regression model. The results showed that the neural network-based model achieved 
higher accuracy in predicting dance performance and assessing the development level of 
physical quality, demonstrating its effectiveness in supporting data-driven dance teaching 
systems.






%%%%%%%%%%%%%%%%%%%%%%%%%%%%%%%%%%%%%%%%%%%%%%%%%%%%%%%%%%%%%%%%%%%%%%%%%%%%%%%%%%%%%%%%%%%%%%%%%%%%%%%%%%%%%%%%%%%%%%%%%%%%%%%%%%%


